Everything's built. Now the moment of truth: we train the model for the
first time and watch what happens.

This chapter is deliberately unheroic. The goal isn't a triumphant
result --- it's to confirm the loop works and to see the model's honest
behavior as it learns. It is the first time a recipe we have been
assembling is finally run: we are not hoping for a masterpiece, we are
checking that the oven works and nothing explodes. And for anyone who
has brewed anything --- beer, coffee, bread --- the \emph{first} batch
is never about the product. It is about learning our equipment, our
ingredients, and our own impatience. This chapter is that first batch.

\section{The data}\label{the-data}

1400 generated tasks: a natural-language specification paired with its
reference Python program. The model's job is to read the spec and
produce the code. Each task is a single semantic record --- one idea ---
rendered three ways (a program, an English instruction, a set of hidden
tests), so the data is internally consistent by construction. We do not
need to worry that a prompt says ``sum'' but its answer subtracts,
because both came from the same seed.

We hold out a separate 300-task test set that shares no operation with
training in a leaky way, so the model cannot simply see the same task
again. This split discipline is not optional: it is what lets us claim,
later, that the model's score reflects unseen tasks, not memorized ones.

\section{The model}\label{the-model}

The student is a small, deliberately plain decoder-only transformer: 4
layers, hidden dimension 128, 4 attention heads, a compact
\textasciitilde291-token vocabulary, and 128-token context. Total
\textbf{884,224 parameters} --- but we report the split honestly,
because it matters more than the round number: \textbf{53,632} are the
embeddings (the token and positional tables, which just store the
vocabulary and position), and \textbf{830,592} are the actual
computation. A tiny model's parameter count is often dominated by its
embedding tables, which are the least interesting part; we never want a
``0.9M parameter model'' claim to hide that most of it is lookup tables.

\section{The training process, step by
step}\label{the-training-process-step-by-step}

Training a small transformer is not a single action; it is a loop
repeated 1,500 times. At each step:

\begin{enumerate}
\def\labelenumi{\arabic{enumi}.}
\tightlist
\item
  We grab a batch of 16 tasks and tokenize each as one sequence: the
  natural-language prompt, a newline, then the reference program.
\item
  The model reads the prompt and predicts each \emph{next} token in the
  combined sequence. This is the standard \textbf{next-token objective}
  --- the model's only job is to guess the next word, given everything
  before it. It is a deceptively simple task, and it is the entire basis
  of how these models learn to write.
\item
  We measure the difference between the model's prediction and the
  actual next token (the cross-entropy loss), take the gradient, and
  nudge the parameters with the \textbf{AdamW} optimizer at learning
  rate \textbf{1e-3}, clipping the gradient magnitude to keep training
  stable.
\item
  Repeat.
\end{enumerate}

We print the loss every 200 steps so we can watch it behave. Here is
what it actually did:

\begin{verbatim}
step    0   loss 5.84
step  200   loss 0.029
step  400   loss 0.022
step  600   loss 0.023
step  800   loss 0.020
step 1000   loss 0.021
step 1200   loss 0.019
step 1499   loss 0.019
\end{verbatim}

That is a large, clean drop --- from 5.84 to 0.019 in 1500 steps. Two
things about this curve are worth sitting with, because they reveal more
than the final number.

\textbf{First, the shape tells us the model is learning, not
memorizing-by-accident.} The loss falls steeply in the first 200 steps
(5.84 → 0.029), then flattens out near a floor. This is a healthy
learning curve: rapid improvement at first, then diminishing returns as
the model fits what it can. If instead the loss had wobbled or stayed
high, we would have had a debugging problem on our hands, not a training
result.

\textbf{Second, the \emph{final} value is not the interesting number.} A
loss of 0.019 means the model can predict the next token of the
\emph{training} sequence almost perfectly. That tells us it fits the
training data --- and nothing more. This is the book's recurring motif,
and it is worth saying loudly: \textbf{a falling training loss, and even
a very low one, is not evidence of learning.} It is evidence of fitting.
The distinction is the entire point of the experiment.

\section{What the training process felt like (the brewing
experience)}\label{what-the-training-process-felt-like-the-brewing-experience}

Anyone who has trained a model knows the peculiar rhythm of it. We set a
batch of runs going and we wait, checking the loss every few minutes
with a combination of hope and suspicion. The first steps are the most
honest --- a fresh, random model has loss near the log of the vocabulary
size, because its best guess is uniform. That initial 5.84 is exactly
what a model that knows nothing looks like.

Then, if things are working, the loss starts to fall, and a strange
thing happens: the drop is \emph{so} fast in the first couple hundred
steps that we wonder whether it is real or whether something is broken.
In this case the model reached near-floor loss within 200 steps. The
rest of the training is just the model quietly polishing the last
fraction of a point.

There is also a genuine practical question that arises during training:
\textbf{how do we know it hasn't crashed mid-run?} In this book we
habitually save the checkpoint \emph{before} evaluating, so a crash
during evaluation costs nothing and a crash during training at least
preserves what was learned up to that point. This is a small discipline,
but it has saved us more than once --- a training run on this hardware
takes real time, and we do not want to lose it to a brittle downstream
step.

\section{The honest caveat}\label{the-honest-caveat}

A falling training loss --- and even a very low one --- is \emph{not}
evidence of learning. We saw this in the pilot and it is the book's
recurring motif. The training loss tells us the model fits the training
data. It does not tell us it can \emph{generalize} to unseen problems,
or that it learned the computation rather than the surface form.

So after training, the only losses that matter are the \emph{held-out}
ones: how well does the model do on problems it never saw? That is
Chapter 10's question, and the answer there is the real measurement of
this brew.

\section{Why the details of the training process
matter}\label{why-the-details-of-the-training-process-matter}

It is tempting to skip the mechanics and jump to results. But the whole
premise of this book is that \emph{the experiment is the subject}, and
an experiment we cannot reproduce is not an experiment. The exact
vocabulary size, the number of layers, the batch size, the learning
rate, the number of steps --- these are not trivia. They are the
conditions under which every later claim is true. If we change one and
the result changes, that tells us something; if we change two and the
result changes, we can no longer say which one caused it. This is the
\textbf{one-variable discipline} that governs every experiment in the
book, and it is why we record the configuration so carefully in a
manifest alongside every result. A reader should be able to take our
numbers, run the same configuration, and get the same outcome --- and if
they cannot, we have not done the work, we have only described it.

\section{What this chapter
establishes}\label{what-this-chapter-establishes}

This first brew is not a result; it is a proof of the loop. We have:

\begin{itemize}
\tightlist
\item
  trained a real model on real data,
\item
  watched its loss descend smoothly,
\item
  saved a checkpoint,
\item
  confirmed that we can take the trained model and produce code for a
  new prompt, run it through the sandbox, and score it with the judge.
\end{itemize}

That loop --- generate data, train, generate code, execute, score --- is
the laboratory working end-to-end. It is the foundation everything else
stands on. The model is built. Now we have to find out what it actually
learned.
